% Generated draft: requires mathematical and source review.
\documentclass[11pt]{article}
\usepackage[margin=25mm]{geometry}
\usepackage{fontspec}
\setmainfont{Latin Modern Roman}
\usepackage{amsmath,amssymb,mathtools,unicode-math}
\setmathfont{Latin Modern Math}
\usepackage{enumitem,needspace,xurl,hyperref}
\setlength{\parindent}{0pt}
\setlength{\parskip}{5pt}
\setlength{\emergencystretch}{3em}
\newcommand{\N}{\mathbb{N}}
\newcommand{\Z}{\mathbb{Z}}
\newcommand{\Q}{\mathbb{Q}}
\newcommand{\R}{\mathbb{R}}
\newcommand{\C}{\mathbb{C}}
\newcommand{\F}{\mathbb{F}}
\newcommand{\A}{\mathbb{A}}
\newcommand{\PP}{\mathbb{P}}
\newcommand{\E}{\mathbb{E}}
\newcommand{\eps}{\varepsilon}
\newcommand{\dd}{\,\mathrm{d}}
\begin{document}
\section*{UnsolvedMath ID: 11000226}
\textbf{Source problem code:} catalogue entry
\subsection{Definitions and mathematical statement}
Let \(S=(X,\omega)\) lie in the intended translation-surface stratum component, and let \(N(S,L)\) be the source's counting function with quadratic asymptotic \(N(S,L)\sim cL^2\). The objects counted, weights, area normalization and dependence of \(c\) are not repeated. Define the error-growth exponent
\[
 \beta(S)=\limsup_{L\to\infty}
    \frac{\log|N(S,L)-cL^2|}{\log L},
\]
with \(\log0=-\infty\) if this extended-real convention is intended.

The source asks:
\begin{enumerate}
\item For a generic translation surface, is \(\beta(S)<2\)? This is a power-saving growth question, stronger than the given \(o(L^2)\) error alone.
\item Is \(\beta(S)\) the same for almost every surface in a fixed connected stratum component? The measure is not printed; normalized area and Masur--Veech measure are the usual interpretation, not an unstated supplied hypothesis.
\item For Veech surfaces in a fixed stratum, is there a uniform bound for the error exponent described by the source as a limit? Ask the same for square-tiled surfaces in that stratum. A square-tiled surface is, up to normalization, assembled from unit squares, giving a branched cover of the square torus. Whether the intended uniform bound must be strictly below \(2\), and whether it may depend on a component, is not explicit.
\end{enumerate}
The source invokes the classical Circle Problem for tori and attributes to Veech [Ve3] a result saying the limsup is actually a limit for Veech surfaces, while saying its value was unknown. This exact-limit claim needs checking against the original counting convention, possible smoothing and error zeros/oscillations; it is not adopted as an unrestricted theorem. No numerical exponent or uniform estimate is supplied here.
\subsection{Short English statement}
Seek a power-saving counting error on generic translation surfaces, almost-sure constancy of its exponent in a stratum component, and uniform exponent bounds for Veech or square-tiled surfaces.
\subsection{Sources}
\begin{itemize}
\item \emph{Problems on Mapping Class Groups and Related Topics}, edited by Farb, Problem 19, ``Error term for counting functions,'' supplied PDF page 263. \url{https://www.math.uchicago.edu/~farb/papers/mcgbook.pdf}. Recorded access 2026-07-29; extraction pdf-font-structure.
\item Veech [Ve3], cited for the limit assertion; full bibliography and precise hypotheses are absent. The classical Circle Problem is named without a separate reference.
\item Local UnsolvedMath registry, \texttt{data/unsolvedmath/problems.json}, permanent ID \(11000226\), exact source code \texttt{catalogue entry}; supplied OPEN-TRIAGE summary reports unverified status.
\end{itemize}
\end{document}
